\subsection{Real PL--TOP comparison in dimensions at most two}

The following comparison is a consequence of the classical joint-family
theorems for three-manifolds and the classical annulus calculation.
The component conventions in those inputs require attention because
the concordance groups here contain all components.

\begin{theorem}\label{thm:low-dimensional-real-concordance-comparison}
Let $B$ be a compact real PL manifold without boundary, every component
of which has dimension at most two.  The natural map
\[
 P_{\mathrm{PL}}(B)\longrightarrow
 \operatorname{Sing}P_{\mathrm{TOP}}(B)
\]
is a weak equivalence.  The source uses jointly PL homeomorphisms on
the whole product with a parameter simplex.  The target is the singular
complex of the compact-open topological concordance group.  In both
categories the bottom is fixed pointwise and the top is free.

If every component of $B$ has dimension zero or one, both concordance
objects are weakly contractible.  For every closed real PL surface $F$
there is also a weak equivalence on the full automorphism groups
\[
 \operatorname{Aut}_{\mathrm{PL}}(F)
 \longrightarrow \operatorname{Sing}\operatorname{Homeo}(F).
\]
Orientable and nonorientable surfaces are included.  The empty base
gives the constant one-element simplicial group.
\end{theorem}

\begin{proof}
For a compact PL manifold $N$, write $\mathcal H_{\mathrm{PL}}(N)$
for the simplicial group of all parameter-preserving joint PL
homeomorphisms of $N\times\Delta^q$, and write
$\mathcal H_{\mathrm{TOP}}(N)=\operatorname{Sing}\operatorname{Homeo}(N)$.
There is no restriction to orientation-preserving maps or to one
component in this notation.

\medskip\noindent\textit{The classical three-manifold comparison.}
Hamstrom allows compact manifolds with boundary and defines her
$\operatorname{Top}(N)$ to be the identity component
\cite[\S0.2, p.~208]{Hamstrom1974}.  Her PL convention is the
semisimplicial complex of joint PL maps on the total parameter
product, with homotopies preserving both the simplex and homotopy
parameters \cite[\S2.1, p.~222]{Hamstrom1974}.  Thus her finite
representatives and relative homotopies compute precisely the
positive homotopy groups of the simplicial group used here.  Her
Theorem~3.1.2, including its proof, gives the natural isomorphisms
\[
 \pi_j\mathcal H_{\mathrm{PL}}(N)
   \xrightarrow{\ \cong\ }
 \pi_j\mathcal H_{\mathrm{TOP}}(N),\qquad j\geq1,\quad\dim N=3,
\]
at the identity \cite[pp.~225--226]{Hamstrom1974}.

Here is the additional component argument.  The Moise--Bing
approximation theorem approximates every homeomorphism of a compact
PL three-manifold by PL homeomorphisms; its all-homeomorphism
formulation appears in
\cite[Chapter~I, \S1, Theorem~1.1, p.~7]{Laudenbach1974}.
Local contractibility of the TOP homeomorphism group
\cite[Corollary~1.1]{EdwardsKirby1971} implies that a sufficiently
close approximation is TOP-isotopic to the original homeomorphism.
Thus every TOP component contains a PL vertex.

For injectivity on components, suppose a TOP isotopy joins two PL
vertices.  Right multiplication by the inverse of its initial vertex
reduces to an isotopy $h_t$ starting at the identity.  Joint PL local
$0$-connectedness gives a $\delta>0$ such that every PL vertex in the
TOP identity component whose uniform distance from the identity is
less than $\delta$ admits a joint PL isotopy to the identity
\cite[Theorem~2.3.1, p.~224, with \S2.1]{Hamstrom1974}.
Subdivide the parameter interval so finely that consecutive values
of $h_t$ have uniform distance less than $\delta/3$.  Approximate
these values by PL maps $g_i$, preserving the two PL endpoints,
with errors less than $\delta/3$ and small enough that each $g_i$
is TOP-isotopic to its corresponding value of $h_t$.  The approximations
therefore belong to the TOP identity component and satisfy
\[
 d(g_{i+1}g_i^{-1},1)=d(g_{i+1},g_i)<\delta.
\]
The local joint PL isotopies of these differences, right-multiplied
by $g_i$, join consecutive approximations.  Their finite concatenation
is jointly PL.  Multiplying back by the initial vertex proves
injectivity on components.  Translation by a PL representative in
each component extends the positive-degree comparison there.  Hence
\begin{equation}\label{eq:three-manifold-full-automorphism-comparison}
 \mathcal H_{\mathrm{PL}}(N)\longrightarrow
 \mathcal H_{\mathrm{TOP}}(N)
 \quad\text{is a weak equivalence when }\dim N=3.
\end{equation}

\medskip\noindent\textit{A product decomposition for surface cylinders.}
First let $F$ be connected and closed, and put $N=F\times I$.
For either category let $\mathcal A(F)\subset\mathcal H(N)$ be the
subgroup preserving the bottom boundary component setwise.  A path
cannot change the permutation of the two boundary components.
Thus $\mathcal A(F)$ is a union of entire components, and
\eqref{eq:three-manifold-full-automorphism-comparison} restricts to
a weak equivalence on these subgroups.

Bottom restriction $r:\mathcal A(F)\to\mathcal H(F)$ has the
product section $s(a)(x,u)=(a(x),u)$.  There is a literal product
decomposition of the underlying simplicial sets,
\begin{equation}\label{eq:bottom-setwise-product-decomposition}
 \begin{aligned}
 \mathcal A(F)&\longrightarrow P(F)\times\mathcal H(F),&
 h&\longmapsto\bigl(hs(rh)^{-1},rh\bigr),\\
 P(F)\times\mathcal H(F)&\longrightarrow\mathcal A(F),&
 (p,a)&\longmapsto ps(a).
 \end{aligned}
\end{equation}
The formulas are inverse and preserve all parameters.  Joint PL
maps have joint PL inverses; in TOP, composition and inversion are
continuous on the compact-domain homeomorphism group.  Thus the TOP
formula is also a product decomposition of the underlying topological
spaces.  The multiplication law is a semidirect product, which does
not affect this underlying product decomposition.

Under \eqref{eq:bottom-setwise-product-decomposition}, the natural
comparison is the product of the concordance comparison and the
automorphism comparison for $F$.  A product of two maps between
nonempty spaces is a weak equivalence only if each factor is one:
bijectivity on the product of component sets implies bijectivity
on each set, and the product decomposition of the based homotopy
groups gives the positive-degree assertion.  Both factors are
nonempty because they contain the identity.  This proves the two
surface assertions for connected $F$.

\medskip\noindent\textit{The circle: positive PL homotopy groups.}
Scott's original definition uses joint PL homeomorphisms on
$\Delta^q\times N$ preserving its parameter projection, with the
usual faces and degeneracies
\cite[Chapter~III, Definition~0.1, p.~19]{Scott1968}.
For the annulus, his natural evaluation map at an interior point
induces isomorphisms
\[
 \pi_j\mathcal H_{\mathrm{PL}}(S^1\times I)
   \longrightarrow\pi_j(S^1\times I),\qquad j\geq1
\]
\cite[Chapter~IV, Theorem~3.1(b)(iii), p.~51,
and its annulus proof, p.~53]{Scott1968}.
Move the evaluation point along a path to the bottom circle.  Evaluating
along this path is a homotopy, with the corresponding basepoint
transport.  Bottom evaluation therefore induces the same
isomorphisms.  Evaluation for $\mathcal H_{\mathrm{PL}}(S^1)$
induces isomorphisms in positive degrees as well: the natural PL
rotation circle is a homotopy equivalence and evaluates with degree
one \cite[Theorem~2.2.4 and proof, pp.~223--224]{Hamstrom1974}.
Since the bottom circle includes into the annulus by a homotopy
equivalence, the evaluation square shows that bottom restriction
induces isomorphisms on every positive homotopy group.
Apply \eqref{eq:bottom-setwise-product-decomposition} to the annulus.
Its first factor $P_{\mathrm{PL}}(S^1)$ has zero positive homotopy
groups.

\medskip\noindent\textit{The circle: components.}
We give the component argument explicitly.  An annulus homeomorphism
fixed on the bottom is orientation preserving.  Indeed, naturality of
the boundary map from its relative fundamental class shows that its
orientation sign equals the degree on each boundary circle; the bottom
degree is one.  Its upper restriction therefore has degree one as well.  Every
orientation-preserving PL circle map has a joint PL isotopy to the
identity.  To construct one, lift it to an increasing PL map
$a:\mathbb R\to\mathbb R$ with $a(x+1)=a(x)+1$, subdivide one
period at all breakpoints, and triangulate each parameter rectangle.
Interpolate the identity on its lower edge and $a$ on its upper
edge affinely on the two triangles.  The spatial slope on either
triangle is respectively $1$ or the positive slope of $a$ on that
interval.  Consequently every horizontal slice is increasing.
The endpoint values respect translation by one, so the finite
joint PL construction descends to the circle.

Set $\rho(s,t)=\min\{s,t\}$.  If $b_s$ is a joint PL isotopy from
the identity to the inverse of the upper restriction, the family
\[
                    (x,t)\longmapsto(b_{\rho(s,t)}(x),t)
\]
is a joint PL concordance isotopy starting at the identity.  Composing
with its endpoint makes the given concordance fix both boundary
circles.  For $\alpha=\{0\}\times I$, the image of $\alpha$ then
has an integer winding $k$ relative to its endpoints in the universal
cover $\mathbb R\times I$.  Composing with
\[
 T_{-k}([x],t)=([x-kt],t)
\]
removes that winding.  This twist lies in the identity component
of the bottom-fixed concordance group, by the joint PL family
$([x],t)\mapsto([x-k\rho(s,t)],t)$.  The upper circle may move
during this family, as the definition permits.

The adjusted arc is homotopic to $\alpha$ relative to its endpoints.
Epstein's relative arc theorem gives a PL ambient isotopy fixed on
both boundary circles which carries the adjusted arc to $\alpha$
\cite[Theorem~3.1 and proof, pp.~89--90]{Epstein1966}.
To correct its parametrization, let $u:I\to I$ be the restriction
on the now setwise fixed arc, and choose a joint PL isotopy $b_s$
from the identity to $u^{-1}$ relative to the endpoints.  In a PL
strip $[-1,1]\times I$ about the arc, the family
\[
 (x,t)\longmapsto\bigl(x,b_{\min\{s,1-|x|\}}(t)\bigr)
\]
fixes the strip boundary and extends by the identity outside it.
At time one its restriction to the central arc is $u^{-1}$.
The formula and its inverse are jointly PL, so postcomposition
makes the arc pointwise fixed.
The adjusted annulus map is still orientation preserving, so it preserves
each of the two local sides of this pointwise fixed arc.  Consequently the
cut-open disk map fixes both copies of the arc pointwise, as well as the
rest of its boundary.
The PL Alexander isotopy of this disk map descends to the annulus
\cite[Theorem~5.2 and proof, pp.~95--96]{Epstein1966}.
This proves $\pi_0P_{\mathrm{PL}}(S^1)=0$ and hence weak
contractibility of the PL concordance group.

\medskip\noindent\textit{The TOP annulus.}
Yagasaki's compact-surface lemma implies that the identity component
of $\operatorname{Homeo}(S^1\times I\ \mathrm{rel}\ S^1\times\{0\})$
is contractible
\cite[Lemma~2.1(ii), p.~2, and proof, p.~3, author version]{Yagasaki2000}.
To check that this is the whole group, take a TOP concordance $h$.
Epstein's approximation isotopy changes it to a PL homeomorphism
$g$ \cite[Theorem~A4 and proof, p.~107]{Epstein1966}; one may choose
any fixed bottom point for the theorem's basepoint.  During that
isotopy $h_t$, the bottom boundary component remains the bottom
component.  If $a_t$ is its bottom restriction, the corrected isotopy
\[
                   (a_t^{-1}\times\operatorname{id}_I)h_t
\]
is fixed on the bottom, starts at $h$, and ends at a PL concordance.
The preceding PL component argument joins that endpoint to the
identity.  Thus the TOP concordance group is connected and equals
its contractible identity component.  The comparison for $S^1$
is now a weak equivalence between weakly contractible objects.

\medskip\noindent\textit{Points and finite disjoint unions.}
For a point, the TOP concordance group is the group of increasing
interval homeomorphisms fixing both endpoints.  Linear interpolation
with the identity contracts it continuously.  In the PL category,
use the joint parameter-preserving Alexander construction
\cite[Theorem~2.2.1 and proof, p.~223]{Hamstrom1974}.
More explicitly, represent a homotopy class by a joint PL map
$h:D\to D$, where $D=\Delta^q\times I$, which preserves the simplex
coordinate and is the identity on $\partial D$.
Let $a$ be the simplex barycenter, put $p=(a,\tfrac12,0)$ in $D\times I$,
and let $C=p*(D\times\{1\})$.  On this cone define
\[
 \widehat h\bigl((1-\lambda)p+\lambda(z,1)\bigr)
   =(1-\lambda)p+\lambda(h(z),1),
 \qquad 0\leq\lambda\leq1,
\]
and extend by the identity outside $C$.  The formula is the identity
on the lateral boundary because $h$ fixes $\partial D$.
Coning a finite triangulation for $h$ makes the extension jointly PL;
coning $h^{-1}$ gives its inverse.  The homotopy coordinate is $\lambda$.
For $z=(u,x)$, the simplex coordinate is
$(1-\lambda)a+\lambda u$ on both sides of the formula, so both
parameters are preserved.  The extension is the identity at time zero
and equals $h$ at time one.  This also applies when $q=0$ and proves
weak contractibility without bilinear interpolation.

A compact manifold has finitely many connected components, and a
bottom-fixed concordance preserves the cylinder of each one.
Both concordance objects are therefore finite products of their
component objects, even when the dimensions differ.  This proves
the stated comparison and weak contractibility for disconnected
bases and includes the empty product.  For a disconnected closed
surface, the same componentwise argument proves the automorphism
comparison on the component-preserving subgroup.  Other components
are indexed by permutations between homeomorphic surface components.
For a homeomorphism between two surface components, apply Moise--Bing
approximation to its product with $\operatorname{id}_I$.  A sufficiently
close PL approximation still carries the bottom boundary to the bottom
boundary; its restriction is a PL homeomorphism of the two surfaces.
Thus all realized permutations have PL representatives.  Translation
by those representatives gives the comparison on the full groups.
\end{proof}

\begin{corollary}\label{cor:effective-low-dimensional-top-models}
Given an ordered finite abstract simplex list promised to realize a
closed PL manifold whose components have dimension at most two,
there is an effective countable simplicial weak model for its real
TOP concordance space.  Its codes have decidable validity, equality
and degeneracy, and computable group and simplicial operations.
Uniformly in the input and the homotopy degree, a complete countable
group presentation can be enumerated; in positive degrees this may
be an integer relation matrix for the abelian homotopy group.
No terminating word problem or manifold-recognition algorithm is
asserted.
\end{corollary}

\begin{proof}
Compose the weak model in
Theorem~\ref{thm:effective-countable-concordance-models} with
Theorem~\ref{thm:low-dimensional-real-concordance-comparison}.
The codes and presentations remain the same; only their
identification with the real TOP homotopy groups is added.
\end{proof}
